\begin{frame}{Los ingresos tributarios aumentaron, pero los ingresos totales disminuyeron}
\centering
\small
\begin{tabular}{lrrr}
\toprule
\textbf{Concepto (\% del PIB)} & \textbf{2024} & \textbf{2025} & \textbf{Cambio} \\
\midrule
Ingresos totales        & 16,4 & 16,3 & -0,1 \\
Ingresos tributarios    & 14,3 & 14,7 &  0,4 \\
Ingresos no tributarios &  0,1 &  0,0 & -0,1 \\
Fondos especiales       &  0,3 &  0,3 &  0,0 \\
Ingresos de capital     &  1,7 &  1,3 & -0,4 \\
\quad Excedentes financieros & 1,5 & 1,2 & -0,3 \\
\quad Reintegros y otros recursos & 0,3 & 0,2 & -0,1 \\
\bottomrule
\end{tabular}
\vspace{0.25cm}
\scriptsize Las diferencias pueden no sumar por redondeo.
\note{[Sources] MFMP 2026, capítulo 1, tabla 1.4.}
\end{frame}

\begin{frame}{El mayor recaudo tributario compensó parcialmente la caída de los ingresos de capital}
\centering
\small
\begin{tabular}{lr}
\toprule
\textbf{Contribución al cambio entre 2024 y 2025} & \textbf{Puntos del PIB} \\
\midrule
Recaudo de impuestos externos &  0,3 \\
Recaudo de renta              &  0,1 \\
Recaudo de IVA                &  0,1 \\
Ingresos no tributarios       & -0,1 \\
Reintegros y otros recursos   & -0,1 \\
Excedentes de Ecopetrol y ANH & -0,3 \\
\midrule
\textbf{Cambio en ingresos totales} & \textbf{-0,1} \\
\bottomrule
\end{tabular}
\note{[Sources] MFMP 2026, capítulo 1, gráfico 1.5.}
\end{frame}